% answers_short.tex —— 2014 年·全国大纲卷·文科（广西） · 简短答案速查
% 由 tools/make_short_answers.py 自 answers.tex 的 \gksolhead/\gkanswer 提取，
% 并与 2014/2014_数学真题_大纲文_答案要点.md 逐题比对后生成，勿手改（改源 MD 后重跑）。

\begin{gknotice}
\textbf{使用说明}：本节只给最终答案便于快速核对，完整推导见解析版。答案为编写者独立推导，\textbf{非官方标准答案}。
\end{gknotice}

\gkpart{选择题}
\begin{center}\begin{tabular}{ccccccccccccc}
\toprule
题号 & 1 & 2 & 3 & 4 & 5 & 6 & 7 & 8 & 9 & 10 & 11 & 12 \\
\midrule
答案 & B & D & C & B & D & B & C & C & A & A & C & D \\
\bottomrule
\end{tabular}\end{center}

\gkpart{填空题}
\begin{itemize}[leftmargin=2.6em]
  \item[\textbf{13.}] $-160$。通项为 $\mathrm C_6^k x^{6-k}(-2)^k$，取 $k=3$ 得系数 $\mathrm C_6^3(-2)^3=-160$。
  \item[\textbf{14.}] $\dfrac32$。$y=\dfrac32-2\left(\sin x-\dfrac12\right)^2\leqslant\dfrac32$，当 $\sin x=\dfrac12$ 时取等。
  \item[\textbf{15.}] $5$。可行域的三个顶点为 $(1,1)$、$(-1,-1)$、$\left(2,\dfrac12\right)$，对应 $z$ 值为 $5$、$-5$、$4$，最大值为 $5$。
  \item[\textbf{16.}] $\dfrac43$。设夹角为 $\theta$，则 $\sin\dfrac{\theta}{2}=\dfrac{\sqrt2}{\sqrt{10}}=\dfrac1{\sqrt5}$，$\tan\dfrac{\theta}{2}=\dfrac12$，故 $\tan\theta=\dfrac43$。
\end{itemize}

\gkpart{解答题}
\begin{itemize}[leftmargin=2.6em]
  \item[\textbf{17.}] (1) 见解析版；(2) $a_n=n^2-2n+2$。由 $b_k=2k-1$ 累加得 $a_n=1+\sum\limits_{k=1}^{n-1}(2k-1)=1+(n-1)^2$
  \item[\textbf{18.}] $B=\dfrac{3\pi}{4}$。由正弦定理得 $3\sin A\cos C=2\sin C\cos A$，即 $3\tan A=2\tan C$，故 $\tan C=\dfrac12$，$A+C=\dfrac{\pi}{4}$。
  \item[\textbf{19.}] (1) 见解析版；(2) $\arccos\dfrac14$。由线面距离 $\sqrt3$ 得 $h^2=3(2-t)^2$，结合 $h^2=4t-t^2$ 及 $D$ 在 $AC$ 上得 $t=1$、$h=\sqrt3$，再算两半平面法向量夹角的余弦为 $\dfrac14$
  \item[\textbf{20.}] (1) $0.31$。$P(X=4)=0.06$，$P(X=3)=0.25$，故 $P(X\geqslant3)=0.31$；(2) $3$。分布列为 $P(X=0)=0.06$、$P(X=1)=0.25$、$P(X=2)=0.38$、$P(X=3)=0.25$、$P(X=4)=0.06$，故 $P(X>2)=0.31\geqslant0.1$，而 $P(X>3)=0.06<0.1$
  \item[\textbf{21.}] (1) $a\geqslant1$ 时在 $\mathbf R$ 上递增；$0<a<1$ 时在 $(-\infty,x_-)$、$(x_+,+\infty)$ 上递增，在 $(x_-,x_+)$ 上递减；$a<0$ 时在 $(x_+,x_-)$ 上递增，在其余两段递减。其中 $x_\pm=\dfrac{-1\pm\sqrt{1-a}}a$；(2) $a\in\left[-\dfrac54,0\right)\cup(0,+\infty)$。由 $f'(x)\geqslant0$ 在 $(1,2)$ 上恒成立得 $a\geqslant-\dfrac{2x+1}{x^2}$ 对 $x\in(1,2)$ 恒成立，右端最大值为 $-\dfrac54$
  \item[\textbf{22.}] (1) $y^2=4x$。由 $|PQ|=\dfrac8p$、$|QF|=\dfrac8p+\dfrac p2$ 及 $|QF|=\dfrac54|PQ|$ 解得 $p=2$；(2) $l:y=\pm(x-1)$。设 $l:x=my+1$，由四点共圆得 $4m-\dfrac4m=0$，故 $m=\pm1$
\end{itemize}
